%% Separately authored additions after exact Many-Valued and Paraconsistent restoration.
% Necessity: the exact inherited chapter does not contain the K3/LP/FDE system
% cards, support-pair matrix, Lukasiewicz and t-norm material, supervaluations,
% or the registered comparison tables, so those attributed additions would
% otherwise vanish from the reader route.
% Definition: this editorial derivative descends from two historical ranges and
% includes current system-card, calculus, theorem, table, and source-role
% repairs; altered ranges are not literal slices.
% Standing: it is editorial reader material, not inherited source or a silent
% proof-status promotion.
% Nonredundancy: the six exact wrappers provide the complete source chapter and
% proof; these slices retain distinct system cards and later topics while the
% historical body and compressed first annex item remain pinned and unreachable.
% Use: parts/p08.tex inputs this file once after SRC-13775-13788; route, slice,
% label, consumer, annex, and publication tests fail closed on any drift.
\noindent\textit{Reader note.} The following system cards, comparison tables,
and later many-valued semantics are separately attributed additions to
the exact inherited chapter above; they neither replace nor alter its claims.

% BEGIN HISTORICAL DERIVATIVE 90-137
\begin{systementry}{entry:p08-k3}{Strong Kleene logic K3}
  \index{Strong Kleene logic}
  \entryaliases{\(\mathbf K_3\); strong three-valued Kleene logic.}
  \entrylanguage{The fixed constant-free propositional language with
    \(\neg,\land,\lor\).}
  \entrysemantics{The matrix with
    \(\mathsf V=\{0,\tfrac12,1\}\), \(D=\{1\}\),
    \(\neg a=1-a\), \(a\land b=\min(a,b)\), and
    \(a\lor b=\max(a,b)\).}
  \entrycalculus{sheet:tableaux-matrix}
  \entryadequacy{proved=annex:p8.k3.1}{For every finite premise set, the
    registered finite-matrix tableau closes exactly for K3 consequence.}
  \entrydecidability{proved=annex:p8.k3.2}{A three-valued truth-table test
    decides finite-premise K3 consequence.}
  \entrycharacteristic{Excluded middle can fail, and there are no tautologies
    in the fixed constant-free language; see
    Theorem~\ref{thm:p08-k3-lp-facts}.}
  \entrytranslations{It is the FDE submatrix on
    \(\{f,n,t\}\), with only \(t\)
    designated.}
  \entryalgebra{The three-element strong Kleene De Morgan algebra with one
    designated value.}
  \entrysources{Kleene1952,OmoriWansing2017}
\end{systementry}

\begin{systementry}{entry:p08-lp}{Logic of Paradox LP}
  \index{Logic of Paradox}
  \entryaliases{\(\mathsf{LP}\).}
  \entrylanguage{The same fixed constant-free language as K3.}
  \entrysemantics{The same value set and strong Kleene operations as K3, but
    with \(D=\{\tfrac12,1\}\).  The middle value is read as carrying both
    positive and negative support.}
  \entrycalculus{sheet:tableaux-matrix}
  \entryadequacy{proved=annex:p8.lp.1}{For every finite premise set, the
    registered finite-matrix tableau closes exactly for LP consequence.}
  \entrydecidability{proved=annex:p8.lp.2}{A three-valued truth-table test
    decides finite-premise LP consequence.}
  \entrycharacteristic{Explosion fails, while the tautologies are exactly the
    classical tautologies in the fixed language; see
    Theorem~\ref{thm:p08-k3-lp-facts}.}
  \entrytranslations{It is the FDE submatrix on
    \(\{f,b,t\}\), with \(b\) and
    \(t\) designated.}
  \entryalgebra{The same three-element strong Kleene algebra as K3 with a
    different designated filter.}
  \entrysources{Priest1979,Priest2006,OmoriWansing2017}
\end{systementry}

% END HISTORICAL DERIVATIVE 90-137

% BEGIN HISTORICAL DERIVATIVE 180-569
\section{First-degree entailment}

\begin{hbdefinition}{def:p08-fde-matrix}{The FDE support-pair matrix}
Let
\[
 \mathrm{BD}=\{(0,1),(0,0),(1,1),(1,0)\}
 =\{f,n,b,t\},
\]
in the displayed order.  For a value \(x=(x^+,x^-)\), the coordinates are its
positive and negative support indicators.  They are needed because a single
linear scale cannot distinguish neither support from both support.  For
\(x=(x^+,x^-)\) and \(y=(y^+,y^-)\), put
\[
\begin{aligned}
 \neg x&=(x^-,x^+),\\
 x\land y&=(x^+\land y^+,\ x^-\lor y^-),\\
 x\lor y&=(x^+\lor y^+,\ x^-\land y^-),
\end{aligned}
\]
where the coordinate operations are Boolean.  Designate
\(\{b,t\}\), exactly the values with positive support.
This finite matrix defines the displayed propositional first-degree-entailment
semantics \cite{Dunn1976,Belnap1977,OmoriWansing2017}.
\index{first-degree entailment}
\index{FDE}
\end{hbdefinition}

\begin{hbdefinition}{def:p08-fde-orders}{Orders on the BD algebra}
For \(x,y\in\mathrm{BD}\), define
\[
\begin{aligned}
 x\leq_t y
 &\quad\text{ iff }\quad
 x^+\leq y^+\ \text{and}\ y^-\leq x^-,\\
 x\leq_k y
 &\quad\text{ iff }\quad
 x^+\leq y^+\ \text{and}\ x^-\leq y^-.
\end{aligned}
\]
The first compares truth support against falsity support; the second compares
total information componentwise.  They are two relations on the same values,
not two consequence relations.  Their nonredundancy is visible at once:
\(n\) and \(b\) are incomparable in the truth order, whereas
they are respectively bottom and top in the information order.  A
\emph{bilattice} is an algebraic structure whose carrier bears two lattice
orders.  The displayed carrier qualifies: truth-order meet and join are
\((\min(x^+,y^+),\max(x^-,y^-))\) and
\((\max(x^+,y^+),\min(x^-,y^-))\), while information-order meet and join take
coordinatewise minima and maxima.  Thus ``bilattice'' records both lattice
orders at once and cannot be replaced by either order alone.  This distinction
is used by the FDE card.
\index{bilattice}
\index{truth order}
\index{information order}
\end{hbdefinition}

\begin{systementry}{entry:p08-fde}{First-degree entailment FDE}
  \index{first-degree entailment!system entry}
  \entryaliases{Belnap--Dunn four-valued logic; \(\mathbf{FDE}\).}
  \entrylanguage{The fixed constant-free propositional language with
    \(\neg,\land,\lor\).}
  \entrysemantics{The support-pair matrix of
    Definition~\ref{def:p08-fde-matrix}, designating
    \(b,t\).}
  \entrycalculus{sheet:tableaux-matrix}
  \entryadequacy{proved=annex:p8.fde.1}{For every finite premise set, the
    registered finite-matrix tableau closes exactly for FDE consequence.}
  \entrydecidability{proved=annex:p8.fde.2}{Finite-premise FDE consequence is
    decidable by exhaustive four-valued evaluation.}
  \entrycharacteristic{It is both paraconsistent and paracomplete: a both-valued
    premise contradiction does not explode, and a neither-valued atom refutes
    excluded middle.}
  \entrytranslations{K3 deletes \(b\) and designates only
    \(t\); LP deletes \(n\) and designates
    \(b,t\).}
  \entryalgebra{\(\mathrm{BD}\) carries the distinct truth and knowledge orders of
    Definition~\ref{def:p08-fde-orders}.}
  \entrysources{Dunn1976,Belnap1977,OmoriWansing2017}
\end{systementry}

\begin{hbtheorem}{thm:p08-k3-lp-facts}{proved=annex:p8.matrix.2}{K3--LP theorem sets and dual consequence}
In the fixed constant-free \(\neg,\land,\lor\)-language:
\begin{enumerate}
  \item \(\mathbf K_3\) has no tautologies;
  \item the LP tautologies are exactly the classical tautologies;
  \item for finite formula sets \(\Gamma,\Delta\), if
        \(\Gamma\entails_S\Delta\) means that every valuation designating all
        members of \(\Gamma\) designates at least one member of \(\Delta\),
        then
        \[
          \Gamma\entails_{\mathbf K_3}\Delta
          \quad\text{ iff }\quad
          \neg\Delta\entails_{\mathsf{LP}}\neg\Gamma,
        \]
        where negation acts memberwise on each finite set.
\end{enumerate}
The third clause is the registered K3--LP duality.  It is not an unqualified
single-conclusion equivalence \cite{Priest1979,OmoriWansing2017}.
\end{hbtheorem}

\begin{hbdefinition}{def:p08-weak-kleene-matrix}{Weak Kleene matrix}
The \emph{weak Kleene matrix} has values
\(\{0,\tfrac12,1\}\), designates only \(1\), fixes
\(\neg\tfrac12=\tfrac12\), and makes \(\tfrac12\) infectious for every
nonunary classical connective: if any immediate argument has that value, so
does the compound; otherwise the operation is classical.  Infectiousness is
an exact operation clause, not the informal claim that a sentence is
meaningless.  This definition is used in Table~\ref{tab:p08-many-valued-matrices}
and nowhere as a substitute for Bochvar's richer external language
\cite{Kleene1952}.
\index{weak Kleene logic}
\end{hbdefinition}

\section{Finite and infinite Lukasiewicz logics}

Finite \L{}ukasiewicz chains use the generic finite-matrix tableau, while the
infinite logic uses its separate implication--negation Hilbert sheet.  The two
calculi are not interchanged: finite-tableau termination does not apply to the
real unit interval, and standard \([0,1]\)-completeness does not identify the
distinct finite theorem sets.

\begin{hbdefinition}{def:p08-lukasiewicz-matrices}{Cignoli's finite chain and H\'ajek's standard \(\text{\L}\) semantics}
For \(n\geq2\), Cignoli, D'Ottaviano, and Mundici write the finite
MV-chain itself as
\[
 \text{\L}_n=\left\{0,\frac1{n-1},\frac2{n-1},\ldots,1\right\}.
\]
The \(n\)-valued Lukasiewicz system uses this value algebra, designated set
\(D=\{1\}\), and the operations
\[
  \neg x=1-x,
  \qquad
  x\Rightarrow y=\min(1,1-x+y).
\]
H\'ajek writes the standard infinite-valued logic as \(\text{\L}\); its
standard value algebra is \([0,1]\).  Its derived strong conjunction is
\(x\mathbin{\ast_{\text{\L}}}y=\max(0,x+y-1)\); weak conjunction and
disjunction are minimum and maximum.  These are H\'ajek's source symbols:
\(\Rightarrow\) is typed here as the truth-function implication and
\(\ast_{\text{\L}}\) as the Lukasiewicz t-norm.  A finite system based on
\(\text{\L}_n\) uses
the generic finite-matrix sheet; the standard \(\text{\L}\) semantics does not
\cite{Lukasiewicz1920,LukasiewiczTarski1930,Hajek1998,
  CignoliDOttavianoMundici2000}.
\index{Lukasiewicz logic}
\index{many-valued logic!Lukasiewicz}
\end{hbdefinition}

\begin{hblemma}{thm:p08-finite-lukasiewicz-tableau}{proved=annex:p8.lukasiewicz.1}{Finite-chain tableau adequacy and decision}
For every fixed \(n\geq2\), instantiate
\hyperref[sheet:tableaux-matrix]{\texttt{sheet:tableaux-matrix}} with
\(\text{\L}_n\), designated set \(\{1\}\), and the operations of
Definition~\ref{def:p08-lukasiewicz-matrices}.  A finite-premise sequent is
valid in that matrix exactly when every corresponding labelled tableau closes;
the exhaustive expansion terminates and decides consequence.
\end{hblemma}

\begin{systementry}{entry:p08-lukasiewicz-finite}{Finite-valued \L{}ukasiewicz logic \(\mathsf L_n\)}
  \entrylanguage{The implication--negation language over the fixed chain
    \(\text{\L}_n\), \(n\geq2\), with other connectives defined by the
    displayed MV operations.}
  \entrysemantics{The finite matrix \(\text{\L}_n\) with designated value
    \(1\), Definition~\ref{def:p08-lukasiewicz-matrices}.}
  \entrycalculus{sheet:tableaux-matrix}
  \entryadequacy{cited=CignoliDOttavianoMundici2000}{Exact specialization in
    Lemma~\ref{thm:p08-finite-lukasiewicz-tableau}.}
  \entrydecidability{cited=CignoliDOttavianoMundici2000}{Finite exhaustive labelled
    expansion decides finite-premise consequence.}
  \entrycharacteristic{\(n=2\) is classical; \(n=3\) is the standard
    three-valued \L{}ukasiewicz matrix; distinct \(n\) need not have identical
    theorem sets.}
  \entrytranslations{The finite chain embeds in the standard MV algebra, but
    the finite tableau is not a calculus for the full interval.}
  \entrysources{Lukasiewicz1920,CignoliDOttavianoMundici2000}
\end{systementry}

\begin{systementry}{entry:p08-lukasiewicz-infinite}{Infinite-valued \L{}ukasiewicz logic \(\mathsf L_{\infty}\)}
  \entrylanguage{Primitive \(\neg,\to\), with standard derived MV
    connectives.}
  \entrysemantics{The standard \([0,1]\) MV algebra with designated value
    \(1\), Definition~\ref{def:p08-lukasiewicz-matrices}.}
  \entrycalculus{sheet:hilbert-lukasiewicz-infinite}
  \entryadequacy{cited=Hajek1998}{Theoremhood and finite-premise standard
    completeness are Theorem~\ref{thm:p08-lukasiewicz-completeness}.}
  \entrydecidability{cited=Hajek1998}{The propositional theoremhood problem is
    decidable; this is not finite truth-table enumeration.}
  \entrycharacteristic{McNaughton functions are exactly the continuous
    piecewise-linear integer-coefficient truth functions.}
  \entrytranslations{Algebraic arbitrary-premise completeness belongs to the
    later MV-algebra entry and is not standard-matrix strong completeness.}
  \entryalgebra{MV algebras.}
  \entrysources{LukasiewiczTarski1930,Hajek1998,McNaughton1951}
\end{systementry}

\begin{landscape}
\begin{longtable}{@{}p{.10\linewidth}p{.16\linewidth}p{.13\linewidth}p{.27\linewidth}p{.25\linewidth}@{}}
\caption{Many-valued matrix comparison.  The theorem-set
column uses the fixed constant-free language for K3, LP, weak Kleene, and FDE.
The finite Lukasiewicz row uses its separately displayed implication language.}
\label{tab:p08-many-valued-matrices}\\
\toprule
System & Values & Designated & Operations & Theorem set and boundary \\
\midrule
\endfirsthead
\toprule
System & Values & Designated & Operations & Theorem set and boundary \\
\midrule
\endhead
\(\mathbf K_3\) & \(0,\tfrac12,1\) & \(1\) & Strong Kleene:
negation \(1-a\), conjunction minimum, disjunction maximum. & Empty theorem
set in the fixed language.  The middle value is nondesignated; finite
multiple-conclusion dual to LP. \\
\addlinespace
\(\mathsf{LP}\) & \(0,\tfrac12,1\) & \(\tfrac12,1\) & The same strong
Kleene operations. & Exactly the classical tautologies, but not classical
consequence: explosion fails.  The middle value is designated. \\
\addlinespace
Weak Kleene matrix & \(0,\tfrac12,1\) & \(1\) &
The middle value is infectious; otherwise operations are classical. & Empty
theorem set in the fixed language.  This is the internal matrix row, not a
card for full Bochvar external logic. \\
\addlinespace
\(\mathbf{FDE}\) & \(f,n,b,t\) &
\(b,t\) & Support-pair De Morgan operations from
Definition~\ref{def:p08-fde-matrix}. & Empty theorem set in the fixed language;
both paraconsistent and paracomplete.  K3 and LP are distinct three-value
submatrices. \\
\addlinespace
The \(n\)-valued \(\text{\L}\) system, with value chain \(\text{\L}_n\),
\(n\geq2\) &
\(0,\frac1{n-1},\ldots,1\) & \(1\) & Lukasiewicz negation and implication
from Definition~\ref{def:p08-lukasiewicz-matrices}. & The finite-valued
Lukasiewicz theorem set.  \(n=2\) is classical; every finite row is decided by
\texttt{sheet:tableaux-matrix}. \\
\bottomrule
\end{longtable}
\end{landscape}

The K3 and LP tableaux share the same expansion shapes because their
operations agree, but their root families differ because their designated sets
differ.  This is the reusable content of the source's first exercise; it does
not identify the two consequence relations.

\begin{hbdefinition}{def:p08-bochvar-external}{Bochvar's external three-valued language}
Extend the weak-Kleene matrix of Definition~\ref{def:p08-weak-kleene-matrix}
by unary assertion operators
\[
 J_T(x)=\begin{cases}1&x=1\\0&x\ne1,\end{cases}\quad
 J_F(x)=\begin{cases}1&x=0\\0&x\ne0,\end{cases}\quad
 J_N(x)=\begin{cases}1&x=\tfrac12\\0&x\ne\tfrac12.\end{cases}
\]
Internal connectives remain infectious; the external operators return only
classical values. The external-language name is necessary because the weak
Kleene internal matrix alone cannot assert that a formula is meaningful,
true, false, or nonsensical. It stands as an expanded truth-functional
language, differs through the three classical-valued assertions, and is used
by the finite-matrix tableau and adequacy statement
\cite{Bochvar1938English1981}.
\index{Bochvar logic}\index{external assertion operator}
\end{hbdefinition}

\begin{hbtheorem}{thm:p08-bochvar-tableau}{proved=annex:p8.bochvar.1}{Finite-matrix adequacy for the Bochvar language}
Instantiate \texttt{sheet:tableaux-matrix} with the three values, designated
set \(\{1\}\), infectious internal operations, and the displayed assertion
tables. A formula is valid exactly when every fully expanded tableau for its
nondesignated root closes. This is ordinary finite truth-table adequacy for the
fixed expanded matrix, not an adequacy claim for a different Bochvar calculus
or for partial valuations.
\end{hbtheorem}

\begin{hbdefinition}{def:p08-post-n-valued}{Post's cyclic \(n\)-valued matrices}
For \(n\ge2\), let \(V_n=\{0,1,\ldots,n-1\}\), designate \(n-1\), put
\(\sim x=x+1\pmod n\), and put \(x\vee y=\max(x,y)\) in the displayed linear
order. The name is needed to fix this cyclic-negation/max family rather than
every logic studied in Post's general classification. It stands as a family of
finite matrices, differs from \L{}ukasiewicz chains because its negation is
cyclic rather than involutive, and is used in the decidability and functional-
generation statement \cite{Post1921}.
\index{Post many-valued logic}\index{cyclic negation}
\end{hbdefinition}

\begin{hblemma}{thm:p08-post-matrix-decision}{proved=annex:p8.post.1}{Decision for each fixed Post matrix}
For each fixed \(n\), validity and finite-premise consequence in the matrix of
Definition~\ref{def:p08-post-n-valued} are decidable by exhaustive valuation.
\end{hblemma}

\begin{hbtheorem}{thm:p08-post-matrix}{cited=Post1921}{Post finite-valued generation boundary}
The displayed Post operations generate the truth functions of the corresponding
finite-valued calculus as established in the cited presentation.  This does
not identify different \(n\), and it does not transfer the two-valued clone
classification of \cite{Post1941} without a separate translation.
\end{hbtheorem}

\begin{hbtheorem}{thm:p08-lukasiewicz-completeness}{cited=Hajek1998}{Lukasiewicz completeness}
For the exact calculus on
\hyperref[sheet:hilbert-lukasiewicz-infinite]{\texttt{sheet:hilbert-lukasiewicz-infinite}}, in the language of
Definition~\ref{def:p08-lukasiewicz-matrices}:
\begin{enumerate}
  \item a formula is a theorem iff every valuation in the standard
        \([0,1]\)-matrix gives it value \(1\);
  \item for a finite premise set, syntactic consequence agrees with preservation
        of value \(1\) in that standard matrix.
\end{enumerate}
These clauses do not license arbitrary-premise strong completeness for the
single standard matrix.  The exact definition of the corresponding algebraic
semantics and its arbitrary-premise completeness theorem are treated in
Section~\ref{sec:p10-many-substructural}; no algebra class is silently introduced
here \cite{CignoliDOttavianoMundici2000}.
\end{hbtheorem}

\begin{hbtheorem}{thm:p08-mcnaughton}{cited=McNaughton1951}{McNaughton's theorem}
A function \(g:[0,1]^m\to[0,1]\) is the truth function of an
\(m\)-variable infinite-valued Lukasiewicz formula iff \(g\) is continuous and
piecewise linear with integer coefficients on every linear piece.  The theorem
does not represent arbitrary continuous functions, and it is not an
arbitrary-premise completeness statement.
\end{hbtheorem}

For each finite \(n\), the finite-matrix theorem in Section~\ref{sec:many-valued-paraconsistent-02} supplies a
terminating decision procedure for consequence in a matrix based on
\(\text{\L}_n\).  The cited
Hilbert axiomatization concerns the infinite logic; no finite-tableau claim is
extended to the real interval, and the deferred algebraic theorem is not used
to strengthen that claim.

\section{T-norm fuzzy logics}

\begin{hbdefinition}{def:p08-tnorm-residuum}{T-norm and residuum}
A \emph{t-norm} is a commutative, associative, coordinatewise monotone
operation \(\ast:[0,1]^2\to[0,1]\) with unit \(1\).  If \(\ast\) is
left-continuous, its \emph{residuum} is
\[
  x\Rightarrow y
  =\sup\{z\in[0,1]:x\ast z\leq y\}.
\]
The t-norm interprets strong conjunction, whereas lattice conjunction remains
minimum; the residuum is a semantic operation interpreting object-language
implication.  The source uses the plain glyph \(\Rightarrow\); its numeric
operands type it as a residuum rather than a sequent arrow.  These distinctions
are needed in every row of Table~\ref{tab:p08-tnorm-logics}
\cite{Hajek1998,CintulaHajekNoguera2011}.
\index{t-norm}
\index{residuum!t-norm}
\end{hbdefinition}

\begin{hbtheorem}{thm:p08-tnorm-completeness}{cited=CintulaHajekNoguera2011}{Propositional t-norm completeness survey}
Let \(\Gamma\cup\{A\}\) lie in the language of the indicated sheet, and let
standard consequence mean preservation of the value \(1\).  If \(\Gamma\) is
finite, then for each
\[
 S\in\{\mathrm{MTL},\mathrm{BL},\mathsf L_{\infty},G,\Pi\},
\]
\(\Gamma\derives_S A\) exactly when every standard \(S\)-evaluation assigning
value \(1\) to every member of \(\Gamma\) assigns value \(1\) to \(A\).
For MTL the standard class ranges over all left-continuous t-norms; for BL it
ranges over all continuous t-norms; the remaining three rows use their single
displayed standard chain.  The restriction that \(\Gamma\) be finite may be
dropped for MTL and \(G\), but not for BL, \(\mathsf L_{\infty}\), or
\(\Pi\).  Arbitrary-premise algebraic completeness uses the corresponding
MTL-, BL-, MV-, G\"odel-, or Product-algebra class, as summarized in the table
and defined in Section~\ref{sec:p10-many-substructural}
\cite{Hajek1998,CignoliEstevaGodoTorrens2000,JeneiMontagna2002}.
\end{hbtheorem}

\begin{landscape}
\begin{longtable}{@{}p{.10\linewidth}p{.18\linewidth}p{.20\linewidth}p{.20\linewidth}p{.23\linewidth}@{}}
\caption{T-norm fuzzy logics.  Every calculus cell links an exact Hilbert
surface; the last column separates standard-chain completeness from
arbitrary-premise algebraic completeness.}
\label{tab:p08-tnorm-logics}\\
\toprule
Logic & Standard t-norm & Residuum & Exact calculus/semantic scope & Completeness surface \\
\midrule
\endfirsthead
\toprule
Logic & Standard t-norm & Residuum & Exact calculus/semantic scope & Completeness surface \\
\midrule
\endhead
\(\mathrm{MTL}\) & Every left-continuous t-norm. &
\(\sup\{z:x\ast z\leq y\}\). &
\hyperref[sheet:hilbert-mtl]{\texttt{sheet:hilbert-mtl}}; standard semantics
ranges over all left-continuous t-norms. & Arbitrary-premise strong standard
completeness over that whole class, and algebraic completeness for MTL-algebras
\cite{EstevaGodo2001,JeneiMontagna2002}. \\
\addlinespace
\(\mathrm{BL}\) & Every continuous t-norm. & The displayed residuum. &
\hyperref[sheet:hilbert-bl]{\texttt{sheet:hilbert-bl}}; standard semantics
ranges over all continuous t-norms. & Finite strong standard completeness over
that class; arbitrary-premise algebraic completeness for BL-algebras
\cite{CignoliEstevaGodoTorrens2000}. \\
\addlinespace
\(\text{\L}\) & \(\max(0,x+y-1)\). &
\(\min(1,1-x+y)\). &
\hyperref[sheet:hilbert-lukasiewicz-infinite]{\texttt{sheet:hilbert-lukasiewicz-infinite}};
the standard operations are fixed in
Definition~\ref{def:p08-lukasiewicz-matrices}. & Finite strong standard
completeness for the standard MV chain; arbitrary-premise algebraic
completeness for all MV-algebras
\cite{Hajek1998,CignoliDOttavianoMundici2000}. \\
\addlinespace
\(G\) & \(\min(x,y)\). & \(1\) if \(x\leq y\), and
\(y\) otherwise. &
\hyperref[sheet:hilbert-godel-fuzzy]{\texttt{sheet:hilbert-godel-fuzzy}};
its standard conjunction satisfies \(x\ast x=x\). & Arbitrary-premise strong
standard completeness for the G\"odel chain, and algebraic completeness for
G\"odel algebras \cite{Hajek1998,CintulaHajekNoguera2011}. \\
\addlinespace
\(\Pi\) & \(x\mathbin{\ast_{\Pi}}y=xy\). & \(1\) if \(x\leq y\), and \(y/x\) otherwise.
& \hyperref[sheet:hilbert-product-fuzzy]{\texttt{sheet:hilbert-product-fuzzy}};
multiplication is its standard t-norm. & Finite strong standard completeness
for the Product chain; arbitrary-premise algebraic completeness for Product
algebras \cite{Hajek1998,CintulaHajekNoguera2011}. \\
\bottomrule
\end{longtable}
\end{landscape}

The table's continuity column is part of the mathematical statement: BL is not
the logic of all left-continuous t-norms, and MTL is not restricted to the
continuous ones.  Likewise, the three named standard chains are examples with
different strong conjunctions, not alternative informal readings of one fixed
matrix.

\section{Supervaluations}

\begin{hbdefinition}{def:p08-supervaluation}{Fine specification spaces and supertruth}
Following Fine's presentation, a \emph{specification space} is a partially
ordered collection of specification-points, where the order records extension
of information.  A \emph{base-point} is the specification-point being
evaluated.  A point above it is a \emph{completion} when it gives a complete
classical specification, and it is \emph{admissible} when it satisfies the
primitive admissibility constraints of that specification space.  The
base-point is \emph{completeable} when at least one complete admissible point
extends it.

At a completeable base-point, a formula is \emph{supertrue} iff it is true at
every complete admissible specification extending that point,
\emph{superfalse} iff it is false at every such specification, and
\emph{indeterminate} otherwise.  Specification-points and admissibility are
metalinguistic semantic data, not matrix truth values.  The extension order is
needed to keep only completions of the base information, while admissibility
can further restrict those completions; neither condition subsumes the other.
This distinction is used in the preservation theorem and example below
\cite{Fine1975}.
\index{supervaluation}
\index{admissible precisification}
\index{supertruth}
\end{hbdefinition}

\begin{hbtheorem}{thm:p08-supervaluation-classical}{cited=Fine1975}{Classical tautologies under supervaluation}
Every classical tautology is supertrue at every completeable base-point.  The
converse holds when every complete classical specification is an admissible
completion of the base-point, but need not hold when primitive admissibility
restricts its completions.
\end{hbtheorem}

\begin{hbexample}{def:p08-supervaluation-gap}{Supertruth with an indeterminate atom}
Suppose two complete admissible specifications extending one base-point
disagree about \(p\).  Then \(p\) is neither supertrue nor superfalse, but every
such specification makes
\(p\lor\neg p\) true, so that compound is supertrue.  Supertruth therefore is
not a truth-functional third value assigned compositionally from the three
supervaluation statuses of immediate constituents.
\end{hbexample}

\section{Sources and navigation for this chapter}

The exact inherited chapter above supplies the matrix definitions, value-labelled
tableau calculus, complete generic theorem and proof, exact K3/LP matrices,
countervaluations, formal-data warning, and exercises at lines 13584--13787.
The K3, LP, and FDE cards all point to the one real generic sheet; their
separately retained finite specializations and decision arguments have distinct
annex proofs.  Finite \L{}ukasiewicz matrices use that same generic sheet,
while the infinite real matrix uses its separate Hilbert sheet.

The new FDE entry follows \cite{Dunn1976,Belnap1977,OmoriWansing2017}; the
Lukasiewicz entry follows
\cite{LukasiewiczTarski1930,McNaughton1951,CignoliDOttavianoMundici2000}; and
the t-norm survey follows
\cite{Hajek1998,EstevaGodo2001,CintulaHajekNoguera2011}.  These citations
support the chapter's matrix, t-norm, notation, and navigation tables.  The Part~8
further-reading section supplies the part-wide source route.
% END HISTORICAL DERIVATIVE 180-569
